\documentclass[11pt]{amsart}

\usepackage[T1]{fontenc}
\usepackage{lmodern}
\usepackage{microtype}
\usepackage{amsmath,amssymb}
\usepackage{booktabs}
\usepackage[hidelinks]{hyperref}

\hypersetup{
  pdftitle={A four-vertex counterexample to Devriendt's Conjecture 7.6}
}

\newtheorem{theorem}{Theorem}
\theoremstyle{remark}
\newtheorem{remark}[theorem]{Remark}

\newcommand{\one}{\mathbf{1}}

\title[A counterexample to Devriendt's Conjecture 7.6]
{A four-vertex counterexample to Devriendt's Conjecture~7.6}
\date{}

\subjclass[2020]{Primary 05C50; Secondary 05C12}
\keywords{effective resistance, resistance curvature, resistance capacity,
submodularity}

\begin{document}

\begin{abstract}
Devriendt conjectured that, for a normalized weighted graph, submodularity of
its resistance capacity implies nonnegative resistance curvature. We give a
normalized weighting of $K_4$ for which all $24$ elementary submodularity
inequalities are strict, while one curvature coordinate equals $-1/126$.
This disproves Conjecture~7.6. The example has minimum possible order and
persists under small perturbations. It does not refute the graph-level
existential statement in Conjecture~7.7.
\end{abstract}

\maketitle

Let $G=(V,E)$ be a finite connected simple graph with $|V|\ge2$ and positive
edge conductances $c$. Let $L$ be its weighted Laplacian and $L^\dagger$ its
Moore--Penrose pseudoinverse. With $e_u$ denoting the standard basis vector at
$u$, define
\[
 \omega_{uv}=(e_u-e_v)^{\mathsf T}L^\dagger(e_u-e_v),
 \qquad
 \Omega=(\omega_{uv})_{u,v\in V},
 \qquad
 K=\Omega^{-1}.
\]
Following \cite[Definitions~3.1, 5.1, and 7.1]{Devriendt}, the weighting is
\emph{normalized} when $\one^{\mathsf T}K\one=1$. For $uv\in E$, put
$r_{uv}=c_{uv}\omega_{uv}$, and define the resistance curvature by
\[
 p_v(c)=1-\frac12\sum_{u:\,uv\in E}r_{uv}.
\]
For normalized $c$, the resistance capacity is the set function
$\tau_\bullet(c)$ given by
\[
 \tau_\varnothing=0,
 \qquad
 \tau_{\{v\}}=\frac12,
 \qquad
 \tau_U=\frac12+\frac12
 \left(\one_U^{\mathsf T}\Omega[U]^{-1}\one_U\right)^{-1}
 \quad (|U|\ge2),
\]
where $\Omega[U]$ is the principal submatrix indexed by $U$ and $\one_U$ is
the all-ones vector in $\mathbb{R}^U$. Devriendt proved that
$p(c)\ge0$ implies submodularity of $\tau_\bullet(c)$ and conjectured the
converse \cite[Theorem~7.5 and Conjecture~7.6]{Devriendt}.

\begin{theorem}\label{thm:main}
Let $G=K_4$ on $V=\{0,1,2,3\}$, order the edges as
$01,02,03,12,13,23$, and set
\[
 c=\frac{3485}{15876}(1,1,2,1,4,4).
\]
Then $c$ is normalized, $\tau_\bullet(c)$ is submodular, and
\[
 p(c)=
 \left(
 \frac{47}{126},
 \frac{20}{63},
 \frac{20}{63},
 -\frac1{126}
 \right).
\]
Moreover, every elementary submodularity slack is positive, with minimum
$1123/69700$. Hence Conjecture~7.6 of \cite{Devriendt} is false.
\end{theorem}

\begin{proof}
Write $w=(1,1,2,1,4,4)$. Its Laplacian is
\[
 L(w)=
 \begin{pmatrix}
 4&-1&-1&-2\\
 -1&6&-1&-4\\
 -1&-1&6&-4\\
 -2&-4&-4&10
 \end{pmatrix}.
\]
Using
\[
 L^\dagger=
 \left(L+\tfrac14\one\one^{\mathsf T}\right)^{-1}
 -\tfrac14\one\one^{\mathsf T},
\]
one obtains
\begin{equation}\label{eq:omega}
 \Omega(w)=\frac1{126}
 \begin{pmatrix}
 0&44&44&35\\
 44&0&36&23\\
 44&36&0&23\\
 35&23&23&0
 \end{pmatrix},
 \qquad
 \one^{\mathsf T}\Omega(w)^{-1}\one=\frac{15876}{3485}.
\end{equation}
Since $\Omega(\lambda w)=\lambda^{-1}\Omega(w)$, the stated choice
$c=(3485/15876)w$ satisfies
$\one^{\mathsf T}\Omega(c)^{-1}\one=1$.

The relative resistances, which are invariant under common rescaling of the
conductances, are
\[
 (r_{01},r_{02},r_{03},r_{12},r_{13},r_{23})
 =
 \left(
 \frac{22}{63},
 \frac{22}{63},
 \frac59,
 \frac27,
 \frac{46}{63},
 \frac{46}{63}
 \right).
\]
Substitution in the definition of $p$ gives the displayed curvature vector.

For distinct $i,j\notin S$, define
\[
 \Delta_{ij}(S)=
 \tau_{S\cup\{i\}}+\tau_{S\cup\{j\}}
 -\tau_S-\tau_{S\cup\{i,j\}}.
\]
Submodularity is equivalent to the nonnegativity of the
$\binom42 2^2=24$ quantities $\Delta_{ij}(S)$. The weighting is invariant
under $1\leftrightarrow2$. Since
$\Omega(c)=(15876/3485)\Omega(w)$, exact inversion of the principal
submatrices in \eqref{eq:omega} gives the following orbit representatives.
In the table, strings such as $23$ denote the corresponding vertex sets.

\begin{center}
\renewcommand{\arraystretch}{1.18}
\begin{tabular}{@{}c@{\qquad}l@{\qquad}l@{}}
\toprule
$ij$ & $S$ & $\Delta_{ij}(S)$ \\
\midrule
$01$ & $\varnothing,2,3,23$
& $\tfrac{713}{6970},\ \tfrac{3888}{17425},\
   \tfrac{2107}{20910},\ \tfrac{2389}{41820}$ \\
$03$ & $\varnothing,1,12$
& $\tfrac{128}{697},\ \tfrac{112}{615},\
   \tfrac{1123}{69700}$ \\
$12$ & $\varnothing,0,3,03$
& $\tfrac{1217}{6970},\ \tfrac{5148}{17425},\
   \tfrac{207}{2788},\ \tfrac{64}{2091}$ \\
$13$ & $\varnothing,0,2,02$
& $\tfrac{1018}{3485},\ \tfrac{3038}{10455},\
   \tfrac{2673}{13940},\ \tfrac{1346}{52275}$ \\
\bottomrule
\end{tabular}
\end{center}

The involution $1\leftrightarrow2$ supplies the four cases for $02$, the four
cases for $23$, and
$\Delta_{03}(\{2\})=\Delta_{03}(\{1\})$, accounting for all $24$
inequalities. Every displayed value is positive, and
\[
 \min_{\substack{i<j\\S\subseteq V\setminus\{i,j\}}}
 \Delta_{ij}(S)
 =\Delta_{03}(\{1,2\})
 =\frac{1123}{69700}>0.
\]
Thus $\tau_\bullet(c)$ is submodular, whereas $p_3(c)<0$.
\end{proof}

\begin{remark}[Minimality and stability]
The connected simple graphs on fewer than four vertices are $K_2$, $P_3$,
and $K_3$. On a weighted tree, each edge has relative resistance $1$, so
$p_v=1-\deg(v)/2\ge0$ on $K_2$ and $P_3$. For a weighted triangle, let
$e,f$ be the edges incident with $v$, let $g$ be the opposite edge, and put
$x_h=c_h^{-1}$ and $X=x_e+x_f+x_g$. The series-parallel law gives
$r_h=1-x_h/X$, and hence
\[
 p_v=1-\frac12(r_e+r_f)=\frac{x_e+x_f}{2X}>0.
\]
Thus order four is minimal. Since curvature and the finitely many elementary
slacks depend continuously on positive conductances, the strict inequalities
$p_3(c)<0$ and $\min\Delta_{ij}(S)>0$ persist under sufficiently small
perturbations, followed by normalization.
\end{remark}

\begin{remark}[Scope]
Conjecture~7.6 is pointwise in the normalized conductance vector. The
displayed Conjecture~7.7 of \cite{Devriendt} is existential over conductance
vectors, so the two formulations have different quantifier structures. The
example above directly refutes Conjecture~7.6, but not the graph-level
statement in Conjecture~7.7: uniform conductances on $K_4$, after
normalization, have curvature $(1/4,1/4,1/4,1/4)$, and hence $K_4$ is
resistance nonnegative.
\end{remark}

\begin{thebibliography}{1}

\bibitem{Devriendt}
K.~Devriendt,
\emph{Graphs with nonnegative resistance curvature},
Ann. Comb. \textbf{30} (2026), 415--438.
\href{https://doi.org/10.1007/s00026-025-00774-x}
{doi:10.1007/s00026-025-00774-x}.

\end{thebibliography}

\end{document}